\documentclass[a4paper,11pt]{article}

\usepackage[T1]{fontenc}
\usepackage[utf8]{inputenc}
\usepackage{lmodern}
\usepackage[margin=2.6cm]{geometry}
\usepackage{amsmath,amssymb,amsthm,mathtools}
\usepackage{booktabs,array}
\usepackage{enumitem}
\usepackage[colorlinks=true,linkcolor=blue,citecolor=blue,urlcolor=blue]{hyperref}
\usepackage[nameinlink]{cleveref}

% Every key equation carries a stable label eq:<slug>. The slugs are referred to
% from the code (docstrings, tests) and are the keys under which formulas are
% exported to the user interface, so do not rename them casually.

\newcommand{\conj}[1]{\overline{#1}}
\newcommand{\T}{T}                       % T_p(z) = (z - p)/(1 - conj(p) z)
\newcommand{\D}{\mathbb{D}}
\newcommand{\bary}{\operatorname{bary}}
\newcommand{\im}{\operatorname{Im}}
\newcommand{\re}{\operatorname{Re}}
\newcommand{\dd}{\mathrm{d}}
\newcommand{\e}{\mathrm{e}}
\newcommand{\ii}{\mathrm{i}}
\newcommand{\code}[1]{\texttt{#1}}

\newtheorem{proposition}{Proposition}
\theoremstyle{remark}
\newtheorem{remark}{Remark}

\title{Kuramoto dynamics on a dumbbell graph:\\ reduction to the hyperbolic disk}
\author{Mathematical notes for the \code{kuramoto-visuals} repository}
\date{Supersedes the formulas of the bachelor thesis \code{v0.pdf}}

\begin{document}
\maketitle

\begin{abstract}
These notes fix the mathematics behind the code in this repository. They start
from the Kuramoto system on a dumbbell graph (two complete graphs joined by a
chain), rewrite it in complex form, reduce each complete-graph component to the
M\"obius group by the Watanabe--Strogatz ansatz, fix the gauge by the conformal
barycenter, and end with the variables and formulas that are actually
integrated and drawn. The starting point is the bachelor thesis kept next to
this file as \code{v0.pdf}; where a formula below differs from the thesis this
is said in a remark. The formulas are backed by numerical tests;
\cref{sec:map} lists which.
\end{abstract}

\tableofcontents

%======================================================================
\section{Setting}\label{sec:setting}

The graph has $2N+l$ vertices: two complete graphs on $N$ vertices and a chain
of $l$ vertices between them. Vertices are numbered
\[
\underbrace{1,\dots,N}_{\text{first clique}},\qquad
\underbrace{N+1,\dots,N+l}_{\text{chain}},\qquad
\underbrace{N+l+1,\dots,2N+l}_{\text{second clique}},
\]
and the chain is attached to vertex $N$ of the first clique and to vertex
$N+l+1$ of the second. We call these two the \emph{attached nodes}; the other
$N-1$ vertices of a clique are its \emph{identical nodes}.

We consider the Kuramoto model \cite{Kuramoto1975} on this graph. All
oscillators have the same natural frequency, which is removed by passing to
a rotating frame, and there is no noise. With symmetric weights $a_{ij}$ the
system is
\begin{equation}\label{eq:full-system}
\dot\theta_i=\sum_{j} a_{ij}\,\sin(\theta_j-\theta_i),\qquad
a_{ij}=\begin{cases}
\lambda/N, & i\neq j \text{ in the same clique},\\
\kappa, & i,j \text{ adjacent on the path } N,\,N+1,\dots,N+l+1,\\
0, & \text{otherwise.}
\end{cases}
\end{equation}
Written out,
\begin{equation}\label{eq:full-system-explicit}
\begin{aligned}
\dot\theta_j&=\frac{\lambda}{N}\sum_{i=1}^{N}\sin(\theta_i-\theta_j), && j=1,\dots,N-1,\\
\dot\theta_N&=\frac{\lambda}{N}\sum_{i=1}^{N-1}\sin(\theta_i-\theta_N)+\kappa\sin(\theta_{N+1}-\theta_N),\\
\dot\theta_j&=\kappa\sin(\theta_{j-1}-\theta_j)+\kappa\sin(\theta_{j+1}-\theta_j), && j=N+1,\dots,N+l,\\
\dot\theta_{N+l+1}&=\kappa\sin(\theta_{N+l}-\theta_{N+l+1})+\frac{\lambda}{N}\sum_{i=N+l+2}^{2N+l}\sin(\theta_i-\theta_{N+l+1}),\\
\dot\theta_j&=\frac{\lambda}{N}\sum_{i=N+l+1}^{2N+l}\sin(\theta_i-\theta_j), && j=N+l+2,\dots,2N+l.
\end{aligned}
\end{equation}
The system is invariant under a common shift of all phases.

For a clique we use its own order parameter, the centroid of its $N$ points on
the unit circle:
\begin{equation}\label{eq:order-parameter}
r=\frac1N\sum_{j\in\text{clique}}\e^{\ii\theta_j},\qquad 0\le |r|\le 1 .
\end{equation}

Both cliques are treated identically, so from \cref{sec:complete} on we write
formulas for the first one, with identical nodes $1,\dots,N-1$, attached node
$N$ and chain neighbor $N+1$. For the second clique replace these by
$N+l+2,\dots,2N+l$, by $N+l+1$ and by $N+l$.

%======================================================================
\section{Complex form}\label{sec:complex}

Put $z_j=\e^{\ii\theta_j}$. A phase equation driven by a complex field $H$,
\begin{equation}\label{eq:riccati}
\dot\theta=\im\bigl(H\,\conj z\bigr)
\qquad\Longleftrightarrow\qquad
\dot z=\tfrac12\bigl(H-\conj H\,z^2\bigr),
\end{equation}
is a Riccati equation for $z$; this follows from $\dot z=\ii z\dot\theta$ and
$\im(H\conj z)=(H\conj z-\conj H z)/2\ii$. Since
$\sum_i a\sin(\theta_i-\theta)=\im\bigl(a\sum_i z_i\,\conj z\bigr)$, the system
\eqref{eq:full-system-explicit} becomes
\begin{align}
\dot z_j&=\frac\lambda2\bigl(r-\conj r\,z_j^2\bigr), \qquad j=1,\dots,N-1,
\label{eq:complex-identical}\\
\dot z_N&=\frac\lambda2\bigl(r-\conj r\,z_N^2\bigr)
+\frac\kappa2\bigl(z_{N+1}-\conj z_{N+1}\,z_N^2\bigr),
\label{eq:complex-attached}\\
\dot z_j&=\frac\kappa2\Bigl[(z_{j-1}+z_{j+1})-\conj{(z_{j-1}+z_{j+1})}\,z_j^2\Bigr],
\qquad j=N+1,\dots,N+l,
\label{eq:complex-chain}
\end{align}
with $r$ the order parameter \eqref{eq:order-parameter} of the first clique,
which includes the attached node $z_N$. In \eqref{eq:complex-identical} the
self-term of the sum cancels, so it may be included.

The point of \eqref{eq:complex-identical} is that the $N-1$ identical nodes obey
the \emph{same} Riccati equation with the same field $H=\lambda r(t)$. Nothing
else about $r$ is used in the reduction: it may depend on time in any way.

%======================================================================
\section{The complete graph}\label{sec:complete}

In this section $\kappa=0$ and we consider one clique alone, so that all $N$
nodes obey \eqref{eq:complex-identical}. This is the classical case
\cite{WatanabeStrogatz1994,MarvelMirolloStrogatz2009,ChenEngelbrechtMirollo2017}.

\subsection{M\"obius action and the Watanabe--Strogatz ansatz}

For $p\in\D=\{|z|<1\}$ let
\begin{equation}\label{eq:mobius}
\T_p(z)=\frac{z-p}{1-\conj p\,z},\qquad \T_p^{-1}=\T_{-p},\qquad \T_p(p)=0 .
\end{equation}
Every conformal automorphism of the disk is $z\mapsto\zeta\,\T_w(z)$ with
$|\zeta|=1$, $w\in\D$; these maps preserve the unit circle and form a
three-parameter group isomorphic to $PSU(1,1)\cong PSL(2,\mathbb R)$, the group
of orientation-preserving isometries of the hyperbolic disk.

A Riccati equation \eqref{eq:riccati} is the infinitesimal form of such a map.
Hence points driven by a common field move by a common automorphism, and one
looks for the solution in the form
\begin{equation}\label{eq:ws-ansatz}
z_j(t)=\zeta(t)\,\T_{w(t)}(\beta_j)=\zeta\,\frac{\beta_j-w}{1-\conj w\,\beta_j},
\qquad |\zeta|=1,\quad |w|<1,\quad |\beta_j|=1,
\end{equation}
with constants $\beta_j$.

\begin{proposition}\label{prop:ws}
Let $H(t)$ be any complex function. The functions \eqref{eq:ws-ansatz} solve
$\dot z_j=\tfrac12(H-\conj H z_j^2)$ for every constant $\beta_j$ on the unit
circle if and only if
\begin{align}
\dot w&=-\tfrac12\,(1-|w|^2)\,\conj\zeta\,H,\label{eq:w-dot-general}\\
\dot\zeta&=-\tfrac12\,\bigl(\conj w\,H-w\,\conj H\,\zeta^2\bigr).\label{eq:zeta-dot-general}
\end{align}
\end{proposition}
The proof is a direct computation, given in Appendix~\ref{app:ws}. With $H=\lambda r$
these are the equations (1.8) of the thesis.

\paragraph{Meaning of the variables.}
\begin{itemize}[nosep]
\item $\zeta$ is a rigid rotation of the whole configuration. We write
      $\zeta=\e^{\ii\alpha}$.
\item $w$ measures how strongly the configuration $\beta$ is squeezed. The map
      $\T_w$ pushes boundary points away from $w/|w|$ and towards $-w/|w|$; as
      $|w|\to1$ every $\beta_j\neq w/|w|$ is sent to $-w/|w|$. Full
      synchronization of the clique is $|w|\to1$.
\item $\beta_j$ are constants of the motion. They are not unique: the triple
      $(w,\zeta,\beta)$ has a three-parameter gauge freedom, fixed in
      \cref{sec:gauge}.
\end{itemize}

\subsection{The closed equation for $w$}

Insert the ansatz into the order parameter:
\begin{equation}\label{eq:q-complete}
q:=\conj\zeta\,r=\frac1N\sum_{j=1}^{N}\T_w(\beta_j),\qquad |q|=|r| .
\end{equation}
Then \eqref{eq:w-dot-general}, \eqref{eq:zeta-dot-general} with $H=\lambda r$
read
\begin{align}
\dot w&=-\frac\lambda2\,(1-|w|^2)\,q,\label{eq:w-dot}\\
\dot\alpha&=-\lambda\,\im\bigl(\conj w\,q\bigr).\label{eq:alpha-dot}
\end{align}
The first equation contains only $w$ and the constants: the $N$-dimensional
dynamics is reduced to a flow on the disk, and $\alpha$ is found afterwards by
a quadrature.

\begin{remark}
In the thesis the $\dot w$ equations (2.8), (2.10), (2.18) are printed with the
denominator $1-w\beta_j$; it is $1-\conj w\beta_j$. The $\dot\zeta$ line of
(2.8) lacks the weights: by \eqref{eq:alpha-dot},
$\dot\zeta=-\tfrac\lambda2\,\zeta\,(\conj w q-w\conj q)$ and
$\conj w\,\T_w(\beta)-\text{c.c.}
=\frac{1-|w|^2}{|1-\conj w\beta|^2}\,(\conj w\beta-w\conj\beta)$.
\end{remark}

\subsection{Gradient flow in the hyperbolic metric}

The Poisson kernel of the disk is
\begin{equation}\label{eq:poisson}
P(w,\beta)=\frac{1-|w|^2}{|w-\beta|^2},\qquad |\beta|=1 .
\end{equation}
It is the factor by which $\T_w$ stretches the circle at $\beta$:
$P(w,\beta)=|\T_w'(\beta)|$. A direct differentiation gives the identity
\begin{equation}\label{eq:poisson-gradient}
(1-|w|^2)\,\frac{\partial}{\partial\conj w}\ln P(w,\beta)=\T_w(\beta).
\end{equation}
Define the potential and the hyperbolic gradient
\begin{equation}\label{eq:potential-complete}
f(w)=\frac1N\sum_{j=1}^{N}\ln P(w,\beta_j),\qquad
\nabla_{\mathbb H}=\tfrac12\,(1-|w|^2)^2\,\frac{\partial}{\partial\conj w},
\end{equation}
the latter being the gradient for the metric
$\dd s^2=4|\dd w|^2/(1-|w|^2)^2$. Summing \eqref{eq:poisson-gradient} over $j$
turns \eqref{eq:w-dot} into
\begin{equation}\label{eq:gradient-flow}
\dot w=-\lambda\,\nabla_{\mathbb H}f(w),\qquad
\frac{\dd f}{\dd t}=-\lambda\,|q|^2=-\lambda\,|r|^2\le0 .
\end{equation}
So $f$ is a Lyapunov function. It tends to $-\infty$ at the boundary away from
the points $\beta_j$, and its only critical point in the disk is where
$\sum_j\T_w(\beta_j)=0$, the conformal barycenter of $\beta$
(\cref{sec:gauge}); this is a maximum, and the flow runs from it to the
boundary.

\begin{remark}
In the thesis the first term of (2.12) is printed with $|w-\beta_j|$; the
Poisson kernel has $|w-\beta_j|^2$.
\end{remark}

\subsection{First integrals}

The cross-ratio
\begin{equation}\label{eq:cross-ratio}
(z_1,z_2,z_3,z_4)=\frac{z_1-z_3}{z_1-z_2}\cdot\frac{z_4-z_2}{z_4-z_3}
\end{equation}
is invariant under every M\"obius map, hence under the flow. Among $N$ points
moved by one automorphism there are $N-3$ independent cross-ratios; these are
the $N-3$ constants of motion of Watanabe and Strogatz.

%======================================================================
\section{The dumbbell: reduction of the cliques}\label{sec:reduction}

Now $\kappa\neq0$. The attached node $N$ obeys \eqref{eq:complex-attached}, not
\eqref{eq:complex-identical}, so it leaves the common M\"obius orbit. The
$N-1$ identical nodes still obey \eqref{eq:complex-identical} with the common
field $\lambda r$, and \cref{prop:ws} applies to them verbatim.

\subsection{The reduced system}

Describe the identical nodes by the ansatz \eqref{eq:ws-ansatz} with
$j=1,\dots,N-1$ and keep the attached node as a variable. The order parameter
of the clique becomes
\begin{equation}\label{eq:q-def}
q:=\conj\zeta\,r=\frac1N\Bigl[\,\sum_{j=1}^{N-1}\T_w(\beta_j)+\conj\zeta\,\e^{\ii\theta_N}\Bigr],
\end{equation}
and the equations are
\begin{align}
\dot w&=-\frac\lambda2\,(1-|w|^2)\,q,\tag{\ref{eq:w-dot}}\\
\dot\alpha&=-\lambda\,\im\bigl(\conj w\,q\bigr),\tag{\ref{eq:alpha-dot}}\\
\dot\theta_N&=\lambda\,\im\bigl(r\,\e^{-\ii\theta_N}\bigr)+\kappa\sin(\theta_{N+1}-\theta_N),
\label{eq:node-dot}\\
\dot\theta_j&=\kappa\sin(\theta_{j-1}-\theta_j)+\kappa\sin(\theta_{j+1}-\theta_j),
\qquad j=N+1,\dots,N+l,\label{eq:chain-dot}
\end{align}
together with the same four equations for the second clique, whose attached
node $N+l+1$ is coupled to $\theta_{N+l}$. In \eqref{eq:node-dot} $r=\zeta q$.
The phases of the identical nodes are recovered from \eqref{eq:ws-ansatz}.

\paragraph{Count.} Each clique contributes $w$ (two real numbers), $\alpha$ and
$\theta_N$; with the chain this is $l+8$ real variables instead of $2N+l$. (The
thesis counts $l+6$, taking each complex $w$ as one.) In each clique the
$N-1$ identical nodes carry $N-4$ independent cross-ratios
\eqref{eq:cross-ratio}, one fewer than for the isolated complete graph.

\subsection{The form with $\psi$}

The thesis keeps the attached node inside the ansatz by letting its ``constant''
move:
\begin{equation}\label{eq:psi-def}
\psi:=\T_w^{-1}\bigl(\conj\zeta\,z_N\bigr),\qquad z_N=\zeta\,\T_w(\psi),\qquad \psi=\e^{\ii\chi}.
\end{equation}
For $\kappa=0$ this is one more constant. In general only the chain term moves
it, divided by $\partial z_N/\partial\psi=\zeta(1-|w|^2)/(1-\conj w\psi)^2$:
\begin{equation}\label{eq:psi-dot}
\dot\psi=-\frac{\kappa}{2\,(1-|w|^2)}
\Bigl[\zeta\,\conj z_{N+1}\,(\psi-w)^2-\conj\zeta\,z_{N+1}\,(1-\conj w\psi)^2\Bigr],
\end{equation}
or, in terms of the angle,
\begin{equation}\label{eq:chi-dot}
\dot\chi=\frac{\kappa\,\sin(\theta_{N+1}-\theta_N)}{P(w,\psi)} .
\end{equation}
Equations \eqref{eq:w-dot}, \eqref{eq:alpha-dot}, \eqref{eq:psi-dot} with
$q=\frac1N\bigl[\sum_{j<N}\T_w(\beta_j)+\T_w(\psi)\bigr]$ are the system
(2.8)--(2.10) of the thesis and are equivalent to
\eqref{eq:w-dot}--\eqref{eq:node-dot}.

The two forms differ in practice. As the clique synchronizes, $|w|\to1$ and
$P(w,\psi)\to0$, so \eqref{eq:chi-dot} is singular: $\psi$ races towards
$w/|w|$ while $\theta_N$ itself moves slowly. The code therefore integrates
$\theta_N$ by \eqref{eq:node-dot} and never forms $\psi$.

\subsection{The potential is no longer a Lyapunov function}

With the moving point $\psi$ the potential \eqref{eq:potential-complete} becomes
\begin{equation}\label{eq:potential}
f(w,\psi)=\frac1N\Bigl[\,\sum_{j=1}^{N-1}\ln P(w,\beta_j)+\ln P(w,\psi)\Bigr],
\end{equation}
and by \eqref{eq:poisson-gradient} the equation for $w$ is still
$\dot w=-\lambda\nabla_{\mathbb H}f$, now for a landscape that moves with
$\psi$. Along solutions
\begin{equation}\label{eq:f-dot}
\frac{\dd f}{\dd t}=-\lambda\,|r|^2
+\frac{2\kappa}{N}\cdot\frac{\im(\conj w\,\psi)}{1-|w|^2}\cdot\sin(\theta_N-\theta_{N+1}).
\end{equation}
The second term has no definite sign and carries the factor $(1-|w|^2)^{-1}$,
so $f$ need not decrease. The derivation is in Appendix~\ref{app:fdot}.

\begin{remark}
In the thesis the second term of (2.16)--(2.17) is printed without the factor
$1/N$, which comes from the $1/N$ in front of $\ln P(w,\psi)$.
\end{remark}

%======================================================================
\section{Gauge and frames}\label{sec:gauge}

\subsection{Gauge freedom}

The ansatz \eqref{eq:ws-ansatz} does not determine $(w,\zeta,\beta)$. For any
automorphism $G$ of the disk, the constants $G(\beta_j)$ together with the map
$\zeta\T_w\circ G^{-1}$ describe the same motion. The thesis and most of the
literature take $w(0)=0$, $\zeta(0)=1$, $\beta_j=z_j(0)$. We use a different,
canonical choice which gives $w$ a meaning independent of the initial instant.

\subsection{Conformal barycenter}

For points $z_1,\dots,z_n$ on the unit circle, $n\ge3$ distinct, the
\emph{conformal barycenter} is the unique $p\in\D$ with
\begin{equation}\label{eq:barycenter}
\sum_{j=1}^{n}\T_p(z_j)=0,\qquad p=:\bary(z_1,\dots,z_n).
\end{equation}
Existence and uniqueness hold whenever no point carries half of the total mass
or more \cite{DouadyEarle1986}. The barycenter is conformally natural: for
every automorphism $M$ of the disk
\begin{equation}\label{eq:bary-equivariance}
\bary\bigl(M(z_1),\dots,M(z_n)\bigr)=M\bigl(\bary(z_1,\dots,z_n)\bigr).
\end{equation}
By \eqref{eq:poisson-gradient} it is the critical point of
$p\mapsto\sum_j\ln P(p,z_j)$.

\paragraph{Computation.} Newton's method in the frame recentered at the current
guess. With $y_j=\T_p(z_j)$, $m=\frac1n\sum y_j$ and $B=\frac1n\sum y_j^2$,
\begin{equation}\label{eq:bary-newton}
d=\frac{m+B\,\conj m}{1-|B|^2},\qquad p\leftarrow\T_{-p}(d)=\frac{d+p}{1+\conj p\,d},
\end{equation}
because $\frac1n\sum\T_d(y_j)=m-d+B\conj d+O(|d|^2)$. If $|d|\ge1$ the step
$d=m$ is taken instead.

\subsection{Canonical gauge}

We require $\bary(\beta_1,\dots,\beta_{N-1})=0$ and $\zeta(0)=1$. Given the
initial phases, with $z_j=z_j(0)$ for the identical nodes,
\begin{equation}\label{eq:calibration}
p_0=\bary(z_1,\dots,z_{N-1}),\qquad \beta_j=\T_{p_0}(z_j),\qquad
w(0)=-p_0,\qquad \zeta(0)=1,
\end{equation}
which reproduces the initial data since $\T_{-p_0}=\T_{p_0}^{-1}$. In this gauge,
by \eqref{eq:bary-equivariance} applied to $M=\zeta\T_w$ and $M(0)=-\zeta w$,
\begin{equation}\label{eq:center}
P:=-\zeta\,w=\bary(z_1,\dots,z_{N-1})\qquad\text{at every time.}
\end{equation}
Thus $P$ is a gauge-independent point of the disk attached to the cluster, its
\emph{center}. It starts wherever the initial configuration puts it and runs to
the boundary as the cluster collapses. Moreover
\begin{equation}\label{eq:centered-frame}
\T_P(z)=\zeta\,\T_{-w}\bigl(\conj\zeta\,z\bigr),\qquad\text{hence}\qquad
\T_P(z_j)=\zeta\,\beta_j ,\quad j=1,\dots,N-1 .
\end{equation}

\subsection{Three frames}\label{sec:frames}

The same state can be drawn in three ways. The user interface shows the first
two, and the variant of the second from \cref{sec:global}; the third is the
frame of the thesis.
\begin{description}[leftmargin=1.5em]
\item[Physical frame.] The oscillators $z_j$ themselves, and the center $P$
      inside the disk.
\item[Centered frame.] Everything mapped by $\T_P$. The cluster center is at
      the origin and by \eqref{eq:centered-frame} the cluster is the fixed
      constellation $\beta$ rotated rigidly by $\zeta$. All other points are
      seen through the same M\"obius lens; when the cluster has collapsed they
      all sit near the single boundary point $\zeta\,\e^{\ii\varphi}$,
      $\varphi=\arg w$.
\item[Gauge frame.] Everything mapped by $z\mapsto\T_w^{-1}(\conj\zeta z)$,
      \begin{equation}\label{eq:gauge-frame}
      \T_w^{-1}\bigl(\conj\zeta\,z_j\bigr)=\beta_j,\qquad
      \T_w^{-1}\bigl(\conj\zeta\,z_N\bigr)=\psi,\qquad
      \T_w^{-1}(0)=w .
      \end{equation}
      The constellation $\beta$ does not move at all, the attached node is the
      point $\psi$ of \eqref{eq:psi-def}, and $w$ is the image of the physical
      origin. This is the frame of the thesis: $w$ moves over the landscape
      $f(\cdot,\psi)$ by \eqref{eq:gradient-flow}.
\end{description}
The gauge frame is the centered frame rotated by $\conj\zeta$.

\paragraph{Potentials in the frames.} Let
$F(p)=\frac1N\sum_{j=1}^{N}\ln P(p,z_j)$ be the potential built from the
current physical points of the clique. Since $P(p,z)=|\T_p'(z)|$ and $\T_P$
inverts $\zeta\T_w$ up to a rotation,
\begin{equation}\label{eq:potential-frames}
F(P)=-f(w,\psi).
\end{equation}
In the physical frame the part of $F$ coming from the identical nodes has its
maximum at $P$ by construction, so that picture shows where the center is, not
a descent. The descent of \eqref{eq:gradient-flow} is seen in the gauge frame.

\subsection{The global center}\label{sec:global}

The centered frame of a cluster is tied to one clique. After that clique has
collapsed it shows the rest of the system as a single boundary point, so it
says nothing about the slow stage. A frame for the whole system is obtained
from the barycenter of \emph{all} nodes,
\begin{equation}\label{eq:global-center}
P_{\mathrm{all}}:=\bary(z_1,\dots,z_{2N+l}),\qquad
\sum_{j=1}^{2N+l}\T_{P_{\mathrm{all}}}(z_j)=0 .
\end{equation}
It exists as long as fewer than half of the nodes coincide; a collapsed cluster
carries the fraction $(N-1)/(2N+l)<\tfrac12$. In the frame $\T_{P_{\mathrm{all}}}$
the configuration is balanced at every time, so it cannot shrink to a point:
when in the physical frame all phases have gathered in a short arc, this frame
shows the shape of that arc magnified. There is no single M\"obius map moving
all nodes, so nothing is rigid in this frame, in contrast to
\eqref{eq:centered-frame}.

The interface also marks the barycenter of the chain nodes alone,
\begin{equation}\label{eq:chain-center}
P_{\mathrm{chain}}:=\bary(z_{N+1},\dots,z_{N+l}),
\end{equation}
next to $P_{\mathrm{all}}$ and the two cluster centers; in a frame centered at
some $p$ each of these points is drawn at its image under $\T_p$.

In double precision $P_{\mathrm{all}}$ is computed from the phases themselves
and is resolved while $1-|P_{\mathrm{all}}|\gtrsim10^{-12}$, that is while the
arc is wider than about $10^{-6}$; for $N=l=20$, $\lambda=1$, $\kappa=10$ this
holds up to $t\approx13\,\tau$.

%======================================================================
\section{Exponential coordinates}\label{sec:coordinates}

Near synchrony $1-|w|$ is exponentially small, so $w$ is a bad variable. Use
the hyperbolic distance $\rho$ from the origin and the direction $\varphi$:
\begin{equation}\label{eq:v-def}
w=\tanh\!\frac\rho2\;\e^{\ii\varphi},\qquad v:=\rho\,\e^{\ii\varphi},\qquad
\rho=2\operatorname{artanh}|w| .
\end{equation}
The complex number $v$ is the coordinate of $w$ in the exponential map of the
hyperbolic disk at the origin; unlike $(\rho,\varphi)$ it is smooth at
$\rho=0$. Equation \eqref{eq:w-dot} becomes
\begin{equation}\label{eq:v-dot}
\dot v=-\lambda\Bigl[\frac{1+c}{2}\,q+\frac{k}{2}\,v^2\,\conj q\Bigr],\qquad
c=\frac{\rho}{\sinh\rho},\qquad k=\frac{1-c}{\rho^2},
\end{equation}
with $c\to1$, $k\to\frac16$ as $\rho\to0$ and $c\to0$ as $\rho\to\infty$. In
components (Appendix~\ref{app:vdot}),
\begin{equation}\label{eq:rho-dot}
\dot\rho=-\lambda\,\re\bigl(\e^{-\ii\varphi}q\bigr),\qquad
\dot\varphi=-\frac{\lambda}{\sinh\rho}\,\im\bigl(\e^{-\ii\varphi}q\bigr).
\end{equation}
The stiff factor $(1-|w|^2)$ is gone and $\rho$ grows at a bounded rate.

\paragraph{State.} The integrated state of a clique is
$(\re v,\im v,\alpha,\theta_N)$, with right-hand side
\eqref{eq:v-dot}, \eqref{eq:alpha-dot}, \eqref{eq:node-dot}.

\paragraph{Late times.} Define the cluster phase and the offset of the attached
node from it,
\begin{equation}\label{eq:cluster-phase}
\Theta:=\alpha+\varphi+\pi=\arg P,\qquad \delta:=\theta_N-\Theta .
\end{equation}
Once the cluster has collapsed, $\T_w(\beta_j)\to-\e^{\ii\varphi}$ and
$\e^{-\ii\varphi}q\to-\frac{N-1}{N}-\frac1N\e^{\ii\delta}$, so that
\begin{equation}\label{eq:rho-late}
\dot\rho\to\lambda\Bigl[1-\frac{1-\cos\delta}{N}\Bigr],\qquad
\dot\Theta\to\frac{\lambda}{N}\sin\delta .
\end{equation}
Hence $\rho=\lambda t+O(1)$, the width of the cluster decays like
$\e^{-\lambda t}$, and the informative quantity is $\rho-\lambda t$, which
converges as $\delta\to0$.

%======================================================================
\section{Numerically stable forms}\label{sec:stable}

In double precision $1-|w|$ is lost at $\rho\approx37$, $\e^{-\rho}$ underflows
at $\rho\approx745$ and $\e^{\rho}$ overflows at $\rho\approx709$, while a run
to full synchrony reaches $\rho$ in the thousands. All formulas below are
written in the small quantity
\begin{equation}\label{eq:s-def}
s:=1-\tanh\frac\rho2=\frac{2\,\e^{-\rho}}{1+\e^{-\rho}},\qquad
1-|w|^2=s\,(2-s),
\end{equation}
which is computed without cancellation and simply becomes $0$ for very large
$\rho$. In particular
\begin{equation}\label{eq:log-one-minus}
\ln(1-|w|^2)=2\Bigl[\ln2-\frac\rho2-\ln\bigl(1+\e^{-\rho}\bigr)\Bigr]
\end{equation}
is finite for every $\rho$.

\paragraph{M\"obius map of a boundary point.} For $\beta=\e^{\ii b}$ put
$\gamma=b-\varphi$, $u=\e^{\ii\gamma}$ and compute
$1-u=-2\ii\sin\frac\gamma2\,\e^{\ii\gamma/2}$. Then
\begin{equation}\label{eq:stable-mobius}
\T_w(\beta)=\e^{\ii\varphi}\,\frac{s-(1-u)}{(1-u)+s\,u}\;\xrightarrow[\;s\to0\;]{}\;-\e^{\ii\varphi}.
\end{equation}

\paragraph{Microscope.} After collapse all phases of the cluster agree to
machine precision, but their arrangement is still determined. Write
\begin{equation}\label{eq:microscope}
\arg z_j-\Theta=\arg\bigl(1-s\,x_j\bigr)=:s\,\mu_j,\qquad
x_j=\frac{1+u_j}{(1-u_j)+s\,u_j},
\end{equation}
which follows from $\T_w(\beta_j)/(-\e^{\ii\varphi})=1-s\,x_j$. For
$s|x_j|<10^{-8}$ one uses $\mu_j=-\im x_j$. The rescaled offsets have a limit,
\begin{equation}\label{eq:microscope-limit}
\mu_j\to-\cot\frac{\gamma_j}{2}\qquad(\rho\to\infty),
\end{equation}
so the collapsed cluster keeps a fixed shape, shrinking like $s$.

\paragraph{Order-parameter deficit.} From
$|r|^2=\frac1{N^2}\sum_{i,j}\cos(\theta_i-\theta_j)$,
\begin{equation}\label{eq:deficit}
1-|r|^2=\frac{4}{N^2}\sum_{i<j}\sin^2\frac{\theta_i-\theta_j}{2},\qquad
1-|r|=\frac{1-|r|^2}{1+|r|},
\end{equation}
where the differences are taken from the microscope:
$\theta_i-\theta_j=s(\mu_i-\mu_j)$ between identical nodes and
$\theta_N-\theta_j=\delta-s\mu_j$. For a collapsed cluster
\begin{equation}\label{eq:deficit-late}
1-|r|^2=\frac{4(N-1)}{N^2}\sin^2\frac\delta2,\qquad
1-|r|\approx\frac{N-1}{2N^2}\,\delta^2 ,
\end{equation}
so at late times $1-|r|$ decays at twice the rate of $\delta$. If $\delta$
changes sign, that is if the attached node passes through the phase of its
collapsed cluster, $1-|r|$ of that clique drops at that instant to the order of
the squared width of the cluster and comes back; this is a property of the
solution, not a numerical effect. The first
formula of \eqref{eq:deficit} applies to any set of phases; the interface also
plots it for the $l$ chain nodes alone, with $N$ replaced by $l$.

\paragraph{Potential.} In \eqref{eq:potential} the point $\psi$ is not formed.
Using $\ln P(w,\psi)=-\ln P\bigl(-w,\T_w(\psi)\bigr)$ and
$|w-\e^{\ii b}|=|(1-u)-s|$,
\begin{equation}\label{eq:stable-potential}
f=\frac1N\Bigl[\sum_{j=1}^{N-1}\bigl(\Lambda-2\ln|(1-u_j)-s|\bigr)
-\Lambda+2\ln|(1-u_N)-s|\Bigr],
\end{equation}
with $\Lambda=\ln(1-|w|^2)$ from \eqref{eq:log-one-minus},
$u_j=\e^{\ii(b_j-\varphi)}$ and $u_N=\e^{\ii(\theta_N-\alpha-\varphi-\pi)}=\e^{\ii\delta}$.
At late times $f\approx\frac{N-2}{N}\Lambda\to-\infty$ linearly in $\rho$.

%======================================================================
\section{Linear analysis and time scales}\label{sec:linear}

\subsection{Energy}

The full system \eqref{eq:full-system} is the gradient flow of
\begin{equation}\label{eq:energy}
E(\theta)=-\sum_{i<j}a_{ij}\cos(\theta_i-\theta_j),\qquad
\dot\theta=-\nabla E,\qquad \dot E=-|\dot\theta|^2\le0 .
\end{equation}
Unlike $f$, the energy is a Lyapunov function of the whole system for all
$\lambda,\kappa$. Its minimum $E=-\sum_{i<j}a_{ij}$ is full synchrony.

\subsection{Synchrony and the slow time scale}

Linearizing \eqref{eq:full-system} at $\theta_1=\dots=\theta_{2N+l}$ gives
$\dot\varepsilon=-L\varepsilon$ with the weighted graph Laplacian
\begin{equation}\label{eq:laplacian}
L=\operatorname{diag}\Bigl(\sum_j a_{ij}\Bigr)-A,\qquad
\varepsilon^{\mathsf T}L\,\varepsilon=\sum_{i<j}a_{ij}(\varepsilon_i-\varepsilon_j)^2\ge0 .
\end{equation}
$L$ is symmetric and positive semidefinite; its kernel is the common phase
shift, and since the graph is connected for $\lambda,\kappa>0$ all other
eigenvalues are positive. Full synchrony is therefore linearly stable. (The
thesis reaches the same conclusion with Gershgorin discs.) The slowest
relaxation is governed by the second eigenvalue,
\begin{equation}\label{eq:tau}
0=\lambda_1<\lambda_2\le\dots,\qquad \tau:=\frac{1}{\lambda_2} .
\end{equation}

There are thus two time scales. The cliques collapse at rate $\lambda$
(\cref{sec:coordinates}); afterwards the chain and the two collapsed clusters
relax at rate $\lambda_2$, which for a long chain is much smaller.
\Cref{tab:tau} compares $\tau$ with the decay time fitted to
$|\theta_N-\theta_{N+l+1}|$ on the window $[2\tau,4\tau]$ of a run of the
reduced system.

\begin{table}[ht]
\centering
\begin{tabular}{rrrrrr}
\toprule
$N$ & $l$ & $\lambda$ & $\kappa$ & $1/\lambda_2$ & fitted $\tau$\\
\midrule
10 & 10 & 1 & 1 & 70.9134 & 70.9137\\
20 & 20 & 1 & 1 & 259.4064 & 259.4127\\
20 & 20 & 5 & 1 & 246.0971 & 246.1038\\
20 & 20 & 1 & 5 & 65.8902 & 65.9081\\
20 & 40 & 1 & 1 & 561.8780 & 561.8788\\
40 & 10 & 2 & 0.5 & 474.0627 & 474.0808\\
\bottomrule
\end{tabular}
\caption{Slow time scale: prediction $1/\lambda_2$ and fit, runs up to $5\tau$,
initial phases uniform with seed $0$. Produced by \code{backend/scripts/timescales.py}.}
\label{tab:tau}
\end{table}

\subsection{The collapsed-cluster system}

The set where the identical nodes of each clique coincide,
$\theta_1=\dots=\theta_{N-1}=\Theta_1$ and
$\theta_{N+l+2}=\dots=\theta_{2N+l}=\Theta_2$, is invariant. On it the system
is exactly
\begin{equation}\label{eq:collapsed}
\begin{aligned}
\dot\Theta_1&=\frac\lambda N\sin(\theta_N-\Theta_1),\\
\dot\theta_N&=\frac{N-1}{N}\,\lambda\sin(\Theta_1-\theta_N)+\kappa\sin(\theta_{N+1}-\theta_N),\\
\dot\theta_j&=\kappa\sin(\theta_{j-1}-\theta_j)+\kappa\sin(\theta_{j+1}-\theta_j),
\qquad j=N+1,\dots,N+l,
\end{aligned}
\end{equation}
and symmetrically at the other end: $l+4$ phases. This is what the dynamics
becomes after the fast stage; the first line is the second relation of
\eqref{eq:rho-late}. Its linearization at synchrony has the same slowest rate
$\lambda_2$ as the full system.

In the phase differences $\Delta_N=\theta_N-\Theta_1$,
$\Delta_j=\theta_j-\theta_{j-1}$ ($N<j\le N+l$) the first equations read
\begin{equation}\label{eq:relative}
\begin{aligned}
\dot\Delta_N&=-\lambda\sin\Delta_N+\kappa\sin\Delta_{N+1},\\
\dot\Delta_{N+1}&=\frac{N-1}{N}\,\lambda\sin\Delta_N-2\kappa\sin\Delta_{N+1}+\kappa\sin\Delta_{N+2},\\
\dot\Delta_j&=\kappa\bigl(\sin\Delta_{j+1}-2\sin\Delta_j+\sin\Delta_{j-1}\bigr),
\qquad N+1<j<N+l .
\end{aligned}
\end{equation}

\begin{remark}
These are the relative variables of section 4.4 of the thesis with the
deviations $\varepsilon_j$ set to zero. In (4.22) of the thesis the coefficient
of $\sin\Delta_N$ in $\dot\Delta_{N+1}$ is printed as $\lambda$; it is
$\frac{N-1}{N}\lambda$.
\end{remark}

\subsection{Stationary states}

On the invariant set of \eqref{eq:collapsed} stationarity forces, from the
first line, $\sin\Delta_N=0$, then from the second $\sin\Delta_{N+1}=0$, and so
on along the chain: every difference is $0$ or $\pi$. A difference equal to
$\pi$ is unstable: the corresponding edges form a cut of the graph on which
all weights enter the Hessian of $E$ as $a_{ij}\cos\pi<0$, and the indicator
vector of one side of the cut is a direction of negative curvature. Hence full
synchrony is the only stable stationary state with collapsed clusters. This
does not classify the stationary states of the full system with non-collapsed
cliques; numerically none other than synchrony is ever reached from uniformly
random initial phases.

\subsection{Two reference states}

Numerically the state at intermediate times is close to ``collapsed clusters
joined by an arithmetic progression along the chain''. Such states are not
stationary; their velocity fields are as follows ($A$ is the step of the
progression).

If the progression runs through the attached nodes and the clusters sit one
step beyond them,
\[
\theta_j=A\cdot\min\{\max\{j,N-1\},N+l+2\}\quad\text{up to a constant,}
\]
then
\begin{equation}\label{eq:state-one}
\dot\theta_j=\sin A\cdot
\begin{cases}
\phantom{-}\lambda/N, & j\le N-1,\\
\phantom{-}\kappa-\frac{N-1}{N}\lambda, & j=N,\\
\phantom{-}0, & N<j\le N+l,\\
-\kappa+\frac{N-1}{N}\lambda, & j=N+l+1,\\
-\lambda/N, & j\ge N+l+2 .
\end{cases}
\end{equation}
If instead the clusters and the attached nodes sit on the ends of the chain,
\[
\theta_j=A\cdot\min\{\max\{j,N+1\},N+l\}\quad\text{up to a constant,}
\]
then
\begin{equation}\label{eq:state-two}
\dot\theta_{N+1}=\kappa\sin A,\qquad \dot\theta_{N+l}=-\kappa\sin A,\qquad
\dot\theta_j=0\ \text{ otherwise.}
\end{equation}
These are (4.10) and (4.13) of the thesis.

%======================================================================
\section{Formulas and code}\label{sec:map}

\Cref{tab:map} lists where each formula is implemented and which test checks
it. Backend modules are in \code{backend/}; ``frontend'' means the file
\code{frontend/src/lib/mobius.ts}; tests are in \code{backend/tests/}.

\begin{table}[ht]
\centering
\footnotesize
\renewcommand{\_}{\textunderscore\allowbreak}
\newcolumntype{L}[1]{>{\raggedright\arraybackslash}p{#1}}
\begin{tabular}{@{}L{0.25\textwidth}L{0.32\textwidth}L{0.38\textwidth}@{}}
\toprule
Formula & Implementation & Test\\
\midrule
\eqref{eq:full-system} & \code{kuramoto.kuramoto\_ode}, \code{build\_dumbbell\_adjacency} & reference for all others\\
\eqref{eq:mobius} & \code{mobius.T}, \code{T\_inv} & \code{test\_inverse\_and\_circle\_preserved}\\
\eqref{eq:ws-ansatz}, \eqref{eq:w-dot}, \eqref{eq:alpha-dot}, \eqref{eq:node-dot}, \eqref{eq:chain-dot}, \eqref{eq:v-dot}
  & \code{reduced.disk\_rhs}, \code{reduced\_rhs}, \code{reconstruct}
  & \code{test\_reduced\_rhs\_is\_the\_full\_rhs}, \code{test\_dumbbell\_matches\_full\_system}\\
\eqref{eq:poisson-gradient}, \eqref{eq:gradient-flow} & \code{mobius.log\_poisson} & \code{test\_gradient\_form\_of\_w\_equation}, \code{test\_isolated\_complete\_graph}\\
\eqref{eq:cross-ratio} & \code{reduced.cross\_ratios} & \code{test\_isolated\_complete\_graph}\\
\eqref{eq:psi-dot}, \eqref{eq:chi-dot} & not used in code & \code{test\_psi\_form\_matches\_theta\_form}\\
\eqref{eq:potential}, \eqref{eq:stable-potential}, \eqref{eq:log-one-minus} & \code{mobius.potential\_gauge}, \code{log\_one\_minus\_w2} & \code{test\_potential\_rate\_identity}\\
\eqref{eq:f-dot} & --- & \code{test\_potential\_rate\_identity}\\
\eqref{eq:barycenter}, \eqref{eq:bary-equivariance}, \eqref{eq:bary-newton} & \code{mobius.barycenter} & \code{test\_barycenter\_zero\_mean\_and\_equivariance}\\
\eqref{eq:calibration} & \code{reduced.calibrate}, \code{initial\_state} & \code{test\_calibration\_reproduces\_initial\_phases}\\
\eqref{eq:center}, \eqref{eq:centered-frame}, \eqref{eq:gauge-frame}, \eqref{eq:potential-frames}
  & frontend \code{clusterCenter}, \code{toFrame}, \code{potential}
  & \code{test\_barycenter\_gauge\_identities}, \code{test\_frames}\\
\eqref{eq:global-center}, \eqref{eq:chain-center} & \code{reduced.global\_frame}, \code{resolved\_barycenter} & \code{test\_dumbbell\_matches\_full\_system}\\
\eqref{eq:v-def}, \eqref{eq:s-def} & \code{mobius.w\_from\_v}, \code{v\_from\_w}, \code{s\_of\_rho} & \code{test\_v\_w\_roundtrip}\\
\eqref{eq:rho-dot}, \eqref{eq:rho-late}, \eqref{eq:deficit-late} & --- & \code{test\_rho\_rate}\\
\eqref{eq:stable-mobius} & \code{mobius.T\_w\_boundary}, frontend \code{tBoundary} & \code{test\_stable\_T\_w\_matches\_naive}, \code{test\_stable\_T\_w\_matches\_mpmath}\\
\eqref{eq:microscope}, \eqref{eq:microscope-limit} & \code{mobius.microscope} & \code{test\_microscope\_definition\_and\_limit}\\
\eqref{eq:deficit} & \code{reduced.order\_deficit}, \code{phase\_deficit} & \code{test\_dumbbell\_matches\_full\_system}\\
\eqref{eq:energy} & \code{kuramoto.energy} & \code{test\_energy\_is\_a\_gradient\_flow}\\
\eqref{eq:laplacian}, \eqref{eq:tau} & \code{kuramoto.fiedler}, \code{reduced.fit\_tau} & \code{test\_linearization\_at\_synchrony\_is\_minus\_laplacian}, \code{test\_dumbbell\_matches\_full\_system}\\
\eqref{eq:collapsed} & --- & \code{test\_collapsed\_cluster\_system}\\
\eqref{eq:state-one}, \eqref{eq:state-two} & --- & \code{test\_near\_observed\_states}\\
\bottomrule
\end{tabular}
\caption{Formula, implementation, test.}
\label{tab:map}
\end{table}

\paragraph{Numerical setup.} The reduced and the full system are integrated
with an explicit Runge--Kutta method of order 8 (\code{DOP853}) at relative
tolerance $10^{-11}$, on output times $t_k\propto\sinh(bk)$ up to $t=3\tau$ by default: the step
is fine during the collapse of the cliques and grows smoothly afterwards. For $N=l=20$, $\lambda=\kappa=1$ and a run up to $t=5\tau$ the
reconstructed phases agree with the full system to about $10^{-8}$ throughout,
while $\rho$ reaches $\approx1300$.

%======================================================================
\appendix

\section{Proof of Proposition~\ref{prop:ws}}\label{app:ws}

Differentiate $z=\zeta(\beta-w)/(1-\conj w\beta)$ at constant $\beta$:
\[
\frac{\dot z}{z}=\frac{\dot\zeta}{\zeta}-\frac{\dot w}{\beta-w}
+\frac{\dot{\conj w}\,\beta}{1-\conj w\beta}.
\]
Let $y=\conj\zeta z=\T_w(\beta)$, so that $\beta=(y+w)/(1+\conj w y)$,
\[
\beta-w=\frac{(1-|w|^2)\,y}{1+\conj w y},\qquad
1-\conj w\beta=\frac{1-|w|^2}{1+\conj w y}.
\]
Substituting and multiplying by $z=\zeta y$,
\[
\dot z=-\frac{\zeta\,\dot w}{1-|w|^2}
+\Bigl[\dot\zeta+\zeta\,\frac{w\,\dot{\conj w}-\conj w\,\dot w}{1-|w|^2}\Bigr]y
+\frac{\zeta\,\dot{\conj w}}{1-|w|^2}\,y^2 .
\]
This must equal $\tfrac12(H-\conj H\zeta^2y^2)$ for all $y$ on the circle, i.e.\
coefficient by coefficient. The constant term gives
\eqref{eq:w-dot-general}; the coefficient of $y^2$ gives its complex
conjugate; the coefficient of $y$ must vanish, which with
\eqref{eq:w-dot-general} is
$\dot\zeta=\zeta(\conj w\dot w-w\dot{\conj w})/(1-|w|^2)
=-\tfrac12(\conj wH-w\conj H\zeta^2)$, that is \eqref{eq:zeta-dot-general}.
\qed

\section{Derivation of \eqref{eq:v-dot}}\label{app:vdot}

Write $q=\e^{\ii\varphi}(q_\parallel+\ii q_\perp)$ with real
$q_\parallel,q_\perp$. From \eqref{eq:w-dot} and $|w|=\tanh\frac\rho2$,
$\dd|w|/\dd\rho=\tfrac12(1-|w|^2)$:
\[
\dot\rho=-\lambda q_\parallel,\qquad
\dot\varphi=-\frac\lambda2\,\frac{1-|w|^2}{|w|}\,q_\perp
=-\frac{\lambda}{\sinh\rho}\,q_\perp ,
\]
which is \eqref{eq:rho-dot}. Then
$\dot v=\e^{\ii\varphi}(\dot\rho+\ii\rho\dot\varphi)
=-\lambda\,\e^{\ii\varphi}(q_\parallel+\ii c\,q_\perp)$. Since
$v^2\conj q=\rho^2\e^{\ii\varphi}(q_\parallel-\ii q_\perp)$,
\[
\frac{1+c}{2}\,q+\frac{1-c}{2\rho^2}\,v^2\conj q
=\e^{\ii\varphi}(q_\parallel+\ii c\,q_\perp),
\]
which gives \eqref{eq:v-dot}.

\section{Derivation of \eqref{eq:f-dot}}\label{app:fdot}

$f$ depends on time through $w$ and through $\chi=\arg\psi$. The first
contribution is, by \eqref{eq:gradient-flow}, $-\lambda|q|^2=-\lambda|r|^2$.
For the second, only the term $\frac1N\ln P(w,\psi)$ depends on $\chi$, and
\[
\frac{\partial}{\partial\chi}\ln P(w,\psi)
=-\frac{\partial}{\partial\chi}\ln|\psi-w|^2
=-\frac{2\,\im(\conj w\,\psi)}{|\psi-w|^2}.
\]
Multiplying by \eqref{eq:chi-dot},
$\dot\chi=\kappa\sin(\theta_{N+1}-\theta_N)\,|\psi-w|^2/(1-|w|^2)$, and by
$1/N$ gives the second term of \eqref{eq:f-dot}.

%======================================================================
\begin{thebibliography}{9}

\bibitem{Kuramoto1975}
Y.~Kuramoto, Self-entrainment of a population of coupled non-linear
oscillators, in \emph{International Symposium on Mathematical Problems in
Theoretical Physics}, Lecture Notes in Physics, Springer, 1975.

\bibitem{WatanabeStrogatz1994}
S.~Watanabe and S.~H.~Strogatz, Constants of motion for superconducting
Josephson arrays, \emph{Physica D} \textbf{74} (1994), 197--253.

\bibitem{MarvelMirolloStrogatz2009}
S.~A.~Marvel, R.~E.~Mirollo and S.~H.~Strogatz, Identical phase oscillators
with global sinusoidal coupling evolve by M\"obius group action, \emph{Chaos}
\textbf{19} (2009), 043104.

\bibitem{ChenEngelbrechtMirollo2017}
B.~Chen, J.~R.~Engelbrecht and R.~Mirollo, Hyperbolic geometry of Kuramoto
oscillator networks, \emph{J.~Phys.~A: Math.~Theor.} \textbf{50} (2017),
355101.

\bibitem{DouadyEarle1986}
A.~Douady and C.~J.~Earle, Conformally natural extension of homeomorphisms of
the circle, \emph{Acta Math.} \textbf{157} (1986), 23--48.

\end{thebibliography}

\end{document}
